% !PRINTED 207
\weberpage{207}
%<<CONT
\[
\xi = \lambda\xi', \quad \eta = \frac{1}{\lambda}\eta',
\]
deren Determinante $= 1$ ist, lässt die Normalform \S.\,65, (4) und in ihr
$a_2'$ ungeändert, während $a_1'$, $a_3'$ in $\lambda^2 a_1'$, $a_3' : \lambda^2$
übergehen, und darin ist $\lambda$ eine willkürliche Grösse.

Wenn wir also
\[
a_2' = z
\]
setzen, so ist
\weq{4}{
I = \varphi(a_1'a_3', z),
}
wenn $\varphi$ eine ganze, rationale, ganzzahlige Function der beiden
Argumente $a_1'a_3'$ und $z$ ist.

Nun ist aber nach \S.\,65, (11)
\[
3a_1'a_3' = z^2 - A,
\]
und wenn wir also (4) mit einer geeigneten Potenz von 3, etwa $3^\nu$,
multipliciren, so erhalten wir
\weq{5}{
3^\nu I = \psi(A, z),
}
worin $\psi$ wieder eine ganzzahlige Function von $A$ und $z$ ist. Nach
\S.\,65, (12) ist aber
\[
z^3 = 3Az - B,
\]
und hiernach können wir alle höheren Potenzen von $z$ durch die erste und
zweite ausdrücken, erhalten also
\weq{6}{
3^\nu I = \chi(A, B, z),
}
worin $\chi$ in Bezug auf $z$ höchstens vom zweiten Grade ist.

Die linke Seite von (6) bleibt aber ungeändert, wenn $z$ durch $a_2'$,
$a_2''$, $a_2'''$, also durch drei verschiedene Werthe, ersetzt wird, und
folglich kann die Function $\chi$ die Variable $z$ überhaupt nicht mehr
enthalten (\S.\,30, II).

Es ist daher
\weq{7}{
3^\nu I = \chi(A, B),
}
worin $\chi$ eine ganzzahlige, ganze, rationale Function ist.

Wenn wir in (7) nach (3)
\[
B^2 = 4A^3 - 27 D
\]
setzen, so folgt, dass (7) eine der beiden folgenden Formen hat
\weq{8}{
3^\nu I = \Phi(A, D), \quad B\Phi(A, D),
}
je nachdem der Grad von $I$ gerade oder ungerade ist.
